% Estilo del PDF del informe: el mismo aspecto que el entregado el 07/10 (serif, A4, índice con puntos).
\usepackage{fontspec}
% Los símbolos de la §5 (↠, ∌, ᵢ, ⁺, ⊆…) se toman de fuentes matemáticas disponibles en TeX Live.
\directlua{luaotfload.add_fallback("simbolos", {"Latin Modern Math:mode=harf;", "DejaVu Sans:mode=harf;"})}
% Mantener las fuentes de la exportación original en Windows; usar equivalentes libres en Linux.
\IfFontExistsTF{Times New Roman}
  {\setmainfont{Times New Roman}[RawFeature={fallback=simbolos}]}
  {\setmainfont{Liberation Serif}[RawFeature={fallback=simbolos}]}
\IfFontExistsTF{Consolas}
  {\setmonofont{Consolas}[Scale=0.9, RawFeature={fallback=simbolos}]}
  {\setmonofont{DejaVu Sans Mono}[Scale=0.9, RawFeature={fallback=simbolos}]}
\usepackage{xcolor}
\definecolor{azul}{HTML}{1F3864}
\definecolor{azulclaro}{HTML}{2E5597}
\definecolor{oro}{HTML}{F2C200}
\usepackage{setspace}
\setstretch{1.15}
% Margen extra para cortar líneas con nombres largos en monoespaciado, que no se pueden partir.
\setlength{\emergencystretch}{3em}
\setlength{\parskip}{4pt}
\usepackage{titlesec}
\titleformat{\section}{\Large\bfseries\color{azul}}{}{0pt}{}
\titleformat{\subsection}{\large\bfseries\color{azul}}{}{0pt}{}
\titleformat{\subsubsection}{\normalsize\bfseries}{}{0pt}{}
\titlespacing*{\section}{0pt}{14pt}{6pt}
\titlespacing*{\subsection}{0pt}{12pt}{4pt}
\titlespacing*{\subsubsection}{0pt}{10pt}{2pt}
% Índice: las secciones en negrita y con línea de puntos, como en el PDF anterior.
\makeatletter
\renewcommand*\l@section[2]{\addpenalty{-\@highpenalty}\vskip 6pt\@dottedtocline{1}{0em}{0em}{\bfseries #1}{\bfseries #2}}
\renewcommand*\l@subsection{\@dottedtocline{2}{1.5em}{0em}}
\makeatother
\usepackage{etoolbox}
\usepackage{needspace}
% Tablas y bloques de dependencias en un cuerpo menor, para que entren en el ancho de la página.
\AtBeginEnvironment{longtable}{\small\setstretch{1.0}}
\makeatletter
\def\verbatim@font{\footnotesize\ttfamily}
\makeatother
\usepackage[font={small,it},labelformat=empty]{caption}
\usepackage{fancyhdr}
\pagestyle{plain}
